\section{Conclusion}
\label{sec:conclusion}

Data selection under distribution shift is a directional problem, and the
criteria the field reaches for by default are magnitude criteria. \method{}
replaces the magnitude question with an alignment question, pays roughly 3
percent of fine-tuning compute to answer it on a sketch, and reaches full-pool
accuracy on four shifted benchmarks with about a third of the tokens. The
excess-risk bound in Proposition~\ref{prop:risk} separates what the curriculum
controls from what it cannot, and the anchor-size corollary it yields is
confirmed by a sharp empirical threshold near 128 examples.

Two limitations bound the result. The method needs a target anchor, so it does
not apply when the deployment distribution is genuinely unknown at training
time. It also assumes the pool contains aligned examples to find; on a pool
with no legal text at all, alignment scoring cannot manufacture any, and our
\textsc{LegalQA-S} gains come from a pool where roughly 4 percent of examples
carry target-relevant structure. Characterising that floor, and combining
alignment with an explicit coverage constraint, are the directions we consider
most likely to pay.

\section*{Acknowledgements}
We thank the Northfield cluster staff for reserved capacity during the sweep,
and the reviewers who pressed us to report the anchor-starvation regime rather
than the anchor sizes that work.
